\chapter{The WebGPU Backend}
\label{ch:webgpu-backend}

\texttt{WEBGPU} is CNA's fifth graphics backend, and its youngest: the project owner
explicitly lifted a former prohibition on this backend and authorized its implementation.
It targets native \texttt{wgpu-native} (pinned to version 29.0.1.1), selected the same way
as every other backend via \texttt{-DCNA\_GRAPHICS\_BACKEND=WEBGPU}. Its own tracking plan
spans fourteen phases and over 130 individually numbered tasks --- large enough that this
chapter covers its trajectory rather than every task.

\section{What ``experimental'' concretely means here}

Unlike the four backends covered in the previous three chapters, \texttt{WEBGPU} is
deliberately excluded from the cross-backend feature matrix used throughout
Chapters~\ref{ch:sdl-renderer}--\ref{ch:vulkan-backend}, on the matrix's own stated grounds
that its feature surface was not yet broad enough for a meaningful side-by-side comparison at
the time that matrix was written. This chapter instead draws on the backend's own dedicated
documentation and milestone history.

\section{The native 2D baseline}

Early milestones established device/surface setup, clear/present, a working
\cnaclass{Texture2D} path, vertex/index buffer uploads, and a WGSL-based
\cnaclass{SpriteBatch} implementation. A 120-frame clear-and-present lifecycle test, and a
manually-reviewed \cnaclass{SpriteBatch} validation scene (covering cropping, tint/alpha,
rotation, flips, filtering, and all three address modes), both passed before the native 2D
baseline was declared closed. The most instructive milestone here, though, is an independent
one: a real, independently-developed XNA-style game (a \emph{Windows Phone Speedy Blupi}
port, from the \texttt{mobile-eggbert} sibling repository) was run against this backend on a
real desktop, reached its main menu, rendered its \cnaclass{SpriteBatch} content
pixel-correctly, and played through a real animated cutscene --- with progress bar and
cross-fade --- triggered by a simulated button click, with no WebGPU validation errors of any
kind. This is a genuinely different, and stronger, kind of proof than any of this backend's
own purpose-built test scenes: an unmodified, independently-authored game simply working.

\section{A bug a manual screenshot review missed, and a pixel assertion caught}

One specific finding is worth naming as a methodological lesson in its own right: the
SpriteBatch pipeline's blend factors did not actually match the non-premultiplied output the
shader itself produced, so a translucent sprite --- one with an alpha value strictly between
fully transparent and fully opaque --- rendered fully opaque instead. This was completely
invisible to the earlier manual screenshot review covering rotation, flips, and filtering; it
was only caught once a real pixel-value assertion replaced ``does this look right in a
screenshot.'' The fix matched Vulkan's own blend-factor pairing. This is the same lesson
Chapter~\ref{ch:spritebatch} drew from SDL\_Renderer's ignored-transform-matrix bug: a
rendering pipeline can look entirely correct in every way a human screenshot review checks,
and still be measurably wrong.

\section{3D: real depth testing, a real gap found while adding it}

The first 3D draw path (\texttt{DrawColoredPrimitives} / \texttt{DrawIndexedColoredPrimitives})
was verified with a genuine near/far depth-ordering test rather than trusting draw order
alone. Building it surfaced a real, previously-invisible gap: applying a
\cnaclass{DepthStencilState} had no effect on this backend at all before this point --- the
apply-state entry point was entirely unimplemented, meaning
\texttt{GraphicsDevice.DepthStencilState} had silently done nothing on WebGPU up to that
point. From there, \cnaclass{PbrEffect} (a real glTF~2.0 metallic-roughness BRDF, ported line
for line from EasyGL's own GLSL implementation) and skinned variants of both
\cnaclass{SkinnedEffect} and \texttt{PbrEffect} (bone-palette skinning up to 72 bones, again
ported directly from EasyGL) were added and both pass their dedicated test suites in full.

\subsection{A worked example: PbrEffect, a NOXNA metallic-roughness material}

\cnaclass{PbrEffect} is worth a concrete look before moving on, since --- unlike every stock
effect in Chapter~\ref{ch:stock-effects} --- it has no real-XNA counterpart at all: real XNA
4.0 predates the glTF~2.0 metallic-roughness PBR model by years, so this is a \texttt{NOXNA}
addition ported line for line from EasyGL's own GLSL implementation, not a port of anything
Microsoft ever shipped. Its property surface mixes ordinary \cnaclass{IEffectMatrices}/%
\cnaclass{IEffectLights} / \cnaclass{IEffectFog} plumbing (identical in shape to every stock
effect) with a genuinely new four-texture-map PBR material:

\begin{lstlisting}
PbrEffect effect(getGraphicsDeviceProperty());
effect.setWorldProperty(worldMatrix);
effect.setViewProperty(camera.View());
effect.setProjectionProperty(camera.Projection());

effect.setTextureProperty(&baseColorMap);               // glTF "baseColorTexture"
effect.setNormalMapProperty(&normalMap);
effect.setMetallicRoughnessMapProperty(&metalRoughMap); // one texture: G=rough, B=metal
effect.setOcclusionMapProperty(&aoMap);
effect.setEmissiveMapProperty(&emissiveMap);

// Both factors multiply the map's channel value (glTF's own "value * factor" convention),
// or stand alone with no map bound:
effect.setMetallicFactorProperty(1.0f);
effect.setRoughnessFactorProperty(1.0f);
effect.setEmissiveFactorProperty(Vector3(0.0f, 0.0f, 0.0f));

effect.setLightingEnabledProperty(true);
effect.EnableDefaultLighting();
\end{lstlisting}

The \texttt{MetallicRoughnessMap} single-texture, two-channel packing above is not an
implementation shortcut CNA introduced --- it is glTF~2.0's own real specified layout for this
material model (green channel roughness, blue channel metallic, red and alpha unused), and
\texttt{PbrEffect} follows it exactly so a metallic-roughness texture exported by any
standards-conformant glTF pipeline can be bound directly with no channel-repacking step. This
is also, concretely, the effect Chapter~\ref{ch:models-meshes}'s two independent skinning
systems both connect to on this backend: the WebGPU port adds a skinned variant of
\texttt{PbrEffect} alongside skinned \cnaclass{SkinnedEffect}, so a glTF character rig
authored with a metallic-roughness material and driven by \cnaclass{AnimationPlayer} reaches
GPU skinning through this effect rather than through \cnaclass{SkinnedEffect} at all.

\section{Render targets, environment mapping, instancing, and mip generation}

\cnaclass{RenderTarget2D} support combines depth and stencil into a single format regardless
of what was requested (the same simplification Vulkan makes, Chapter~\ref{ch:vulkan-backend}),
and needed a real type-safety fix when sampling a render target back as an ordinary texture.
MSAA support, once implemented, initially appeared to fail its own dedicated test --- but
investigation found the infrastructure was already correct, and the reported failure was a
defect in the \emph{test itself} (a missing \texttt{RasterizerState::CullNone}), not in the
MSAA implementation. \cnaclass{EnvironmentMapEffect} and genuine hardware instancing were
added together, with instancing implemented as a real per-instance vertex stream needing no
bind-group changes, unlike the effect families that do require one per draw.
\cnaclass{RenderTargetCube} support surfaced a new, cross-cutting, still-unfixed bug worth
flagging directly: sprite geometry is computed from the back buffer's own logical dimensions
unconditionally, never from whichever render target is currently bound --- meaning a
\cnaclass{SpriteBatch::Draw()} call into a smaller or differently-sized off-screen target may
not actually cover it the way a caller would expect. Mip generation for \cnaclass{Texture2D}
and \cnaclass{TextureCube} is implemented as a genuine filtered downsample (a series of
full-screen-triangle render passes, since \texttt{wgpu-native} v29 has no direct
blit-image-equivalent call) --- and, notably, this backend auto-regenerates mip content on
every level-0 write to a plain \cnaclass{Texture2D} / \cnaclass{TextureCube}, which neither real
XNA/FNA nor any other CNA backend does (every other backend only regenerates mips on
render-target unbind). This is documented explicitly as a deliberate divergence that makes
this specific behavior strictly \emph{better} than every other backend's, while still being
genuinely different --- worth knowing if you are relying on mip regeneration timing to match
across backends.

\subsection{A worked example: the still-open render-target-size trap}

The \cnaclass{RenderTargetCube} sprite-geometry bug named above is easy to reproduce
mentally, and worth walking through precisely because it produces no error of any kind ---
only a subtly wrong result:

\begin{lstlisting}
// e.g. a 256x256 render target, while the game window itself is 1920x1080:
device.SetRenderTarget(&smallOffscreenTarget);
spriteBatch.Begin();
spriteBatch.Draw(sourceTexture, Vector2::Zero, Color::White); // "fill the target"
spriteBatch.End();
\end{lstlisting}

On a backend without this bug, a sprite meant to cover the entire currently-bound render
target is sized and positioned relative to \emph{that target's own} dimensions. On this
backend, sprite geometry is instead computed from the back buffer's logical dimensions
unconditionally --- so the same \texttt{Draw()} call above, intended to fill a
$256\times256$ off-screen target, is actually sized as though it were filling the full
$1920\times1080$ window, silently covering far more (or, depending on aspect ratio,
differently-shaped) area than the render target it is actually being drawn into. Nothing about
this throws or logs; the only symptom is visually wrong content in the resulting render
target, which then propagates into whatever samples it afterward. The practical mitigation
until this is fixed: avoid drawing \cnaclass{SpriteBatch} content into a render target whose
dimensions differ from the back buffer's own on this specific backend, or verify the result
with an actual pixel readback rather than trusting the call to behave as documented.

\section{What remains open}

Several gaps are shared with every other CNA backend rather than unique to this one: no
backend anywhere in the project does real block-compressed texture upload, and
\texttt{Texture2D::GetData()} readback from an arbitrary texture remains unimplemented here
too. Specific to this backend: multiple simultaneous render targets, the stencil half of
\cnaclass{DepthStencilState} (the depth half is done), full \cnaclass{BlendState}/%
\cnaclass{RasterizerState} culling and wireframe mapping, viewport/scissor state, custom
\cnaclass{SpriteBatch} effects, and custom WGSL effects authored outside the stock-effect
set. Browser/Emscripten WebGPU is entirely out of scope so far --- this backend is native
\texttt{wgpu-native} only.

\section{Verification methodology, and how it differs from D3D9's}

This backend's own verification runs on a native automated CTest suite (gracefully skipping
where no Wayland/X11 display is available) plus a battery of dedicated pixel-assertion tests
per feature area, and the one independent real-game integration test described above. It is
explicitly \emph{not} covered by the oracle-versus-real-XNA methodology
Chapter~\ref{ch:direct3d-backends} describes for \texttt{D3D9} / \texttt{D3D11} / \texttt{D3D12}
--- that comparison against a genuine, running XNA~4.0 reference is reserved for the three
Direct3D backends specifically, not applied here.
